\section{习题提示与解答要点}\label{app:exercises}

各讲章末已给出练习与部分行内提示。本附录只解决“卡住时看哪里“：按讲编号给出
思路起点、常见错法与可自查的数值。[答] 标的是可直接核对的数值结论。
所有超出原文的推导都标 \ours，与正文标记一致。

\subsection*{第 \ref{sec:cs} 讲（Chern--Simons 作用量）}
\begin{enumerate}[leftmargin=*]
  \item 变分：先分部积分把 $\int\Tr(A\wedge\dd\,\delta A)$ 写成 $-\int\Tr(\dd A\wedge\delta A)$，
        再用 $\delta(A\wedge A\wedge A)=3\,\delta A\wedge A\wedge A$（在迹下循环）。
        常见错法：漏掉 $[A_i,A_j]$ 项的因子 $2/3\times3=2$。
  \item 取 $g:S^3\to SU(2)$ 为一次绕数映射，则 $m=\frac{1}{24\pi^2}\int\Tr(g^{-1}\dd g)^3=1$，
        代入 (D-2.2) 得 $\Delta L=2\pi k$。\hfill[答 $m=1$]
  \item $k\in\mathbb Z+\frac12$ 时 $\exp(\ii L)\to-\exp(\ii L)$，单值性失效。
        这说明“整性“是\textbf{对场配置空间整体}的要求，而不是逐点条件；
        若允许自旋结构（标架），可松动为 $k$ 与标架反常的配合（第 \ref{sec:surgery} 讲）。
  \item $SO(N)$ 的 $N$ 维表示是实表示，其 Wilson 线的量子化给出 Kauffman 多项式；
        $SU(N)$ 定义表示给 Jones。关键不在群大小，而在表示是否自共轭。
  \item 平稳相位：$k\to\infty$ 时 (1.6) 被平坦联络（$F=0$）主导，即 §\ref{sec:phase} 的 $\cM$；
        与 (3.5) 的大 $k$ 极限 $\Inv(\otimes_iR_i)$ 对照，可看出“大 $k$ = 经典“。
\end{enumerate}

\subsection*{第 \ref{sec:phase} 讲（相空间与辛商）}
\begin{enumerate}[leftmargin=*]
  \item 复算 D-3.1：把 $L$ 写成 $\int\dd t\int_\Sigma(\cdots)$，注意 $\varepsilon^{0ij}$ 只留两个空间指标，
        三阶项自动消失。检查点：时间一阶导的系数给出 (3.2)。
  \item $G=SU(2)$、$d=3$、$g=2$：$\dim_{\mathbb R}\cM=(2g-2)d=6$，复维数 $3$。
        \hfill[答 $6$ / $3$]
  \item 两个变元同时规范变换：用 $\omega$ 的规范不变性 + 高斯约束消去纵向分量；
        关键是把 $\delta A$ 分解为纵向（规范）与横向（物理）两部分。
  \item 归一化 $\hat G=\frac{k}{4\pi}G$：验证 $\hat G$ 与 $\omega$ 的关系在 $k\to k+h^\vee$ 下的平移。
  \item 自由度计数：约束 $d$ 个、规范 $d$ 个，每点相减为 $0$；但整体（拓扑）部分留下 $2g$ 个
        模，这就是 $\dim\cM=(2g-2)d$ 的来源（$g=1$ 时退化为 $0$）。
  \item 紧 + 有限辛体积 $\Rightarrow$ 有限维：这是半经典计数 $\mathrm{Vol}/(2\pi)^n n!$ 的前提；
        不是严格证明（严格版本要指标定理）。
  \item 错在：$g=1$ 时 $\dim\cM=0$ 但 $\cM$ 可以有多个连通分支（离散的平坦联络），
        态空间维数是分支数（可积表示数），不是 $1$。$SU(2)$ 给 $k+1$。
\end{enumerate}

\subsection*{第 \ref{sec:prequant} 讲（预量子化与极化）}
\begin{enumerate}[leftmargin=*]
  \item 用 $[\nabla_X,\nabla_Y]s=\nabla_{[X,Y]}s-\ii\hbar\omega(X,Y)s$ 定义预量子化联络；
        注意这里的 $\hbar$ 与本文 $k$ 的换算见附录 \ref{app:conventions}。
  \item 由陈类整性 $\frac{1}{2\pi}[\omega]\in H^2(M;\mathbb Z)$ 与 $\omega_{\rm pre}=\omega/\hbar$，
        得 $\frac{1}{2\pi\hbar}[\omega]\in H^2(M;\mathbb Z)$。\ours
  \item $\mathbb{CP}^1$ 面积 $A$：可预量子化 $\iff A/(2\pi\hbar)\in\mathbb Z$。
  \item $\nabla^{0,1}$ 平方为零 $\iff F^{0,2}=0$ $\iff$ 全纯结构存在；Kähler 时自动成立。
  \item $K_{G/T}=\cL_{-2\rho}$：用 $T^*(G/T)\cong G\times_T\mathfrak n$（$\mathfrak n$ 为幂零部分），
        其行列式线丛的权是 $2\rho$；故 $K^{1/2}=\cL_{-\rho}$。
  \item 错在：本文不写整性，是因为行列式线丛由构造给出（$c_1$ 恒整），不是不需要。
  \item 半形式：$(K^{1/2})$ 的加入把 $k$ 平移 $h^\vee$，但本文全文只用 $\cL^{\otimes k}$，
        所以文中 $k$ 与带半形式的 $k$ 差 $h^\vee$——这是约定，不是错误。
  \item $\omega\to-\omega$：预量子化联络变共轭、$q\to q^{-1}$、$S\to S^*$（至少两处）。
\end{enumerate}

\subsection*{第 \ref{sec:blocks} 讲（共形块）}
\begin{enumerate}[leftmargin=*]
  \item p.370 三句：断言（$J$ 无关 $\iff$ 平坦联络）、引证（Segal 证平坦）、限定（只有射影平坦）。
        删掉限定会把“射影平坦丛“误读成“平坦丛“，从而把 $\cH_\Sigma$ 当成严格函子。
  \item 射影平坦 $\iff$ 曲率 $=\omega\cdot\mathrm{id}$；单值性给射影表示，反常给因子的量子化。
  \item KZ $n=2$：解为 $(z_1-z_2)^{-2h_j}$，指数由 $h_j=\frac{j(j+1)}{k+2}$ 给。\hfill[答 $h_{1/2}=3/8$]
  \item 判据：$g=0,n=3$ 与 $g=1,n=0$ 平坦（有典范联络）；$g=0,n=5$、$g=2,n=0$ 只射影平坦。
  \item $c=k\dim G/(k+h^\vee)$：$SU(2),k=1$ 给 $c=1$；$SU(3),k=2$ 给 $c=\frac{2\cdot8}{2+3}=\frac{16}{5}$。
  \item (M1)--(M6) 中只有 (M1)(M2) 有原文断言；函子级对应在 p.397 只有 “obvious”。
  \item 错在：把“单值性非平凡“直接推成“非平坦“；单值性只说明不能平凡化，射影平坦仍可平坦化到射影丛。
  \item $\cH_{-\Sigma}\cong\cH_\Sigma^*$ 到射影级；射影因子在定向反转下变号（p.370 原话）。
\end{enumerate}

\subsection*{第 \ref{sec:wilson} 讲（Wilson 线与 BWB）}
\begin{enumerate}[leftmargin=*]
  \item 源项：对 $A_0$ 变分得 (3.4)，右端是 $\bigotimes_s\End(V_{R_s})$ 中的算符，故不能与经典 $F=0$ 并列。
  \item $U(1)$ 单电荷：解给 $A\sim\frac{e}{2\pi}\dd\theta$（绕 $P$ 的角），绕数即电荷；奇异性是必然的。
  \item Weyl 维数：$SU(2)$ 的 $R_j$ 维数 $2j+1$；$SU(3)$ 基本表示 $3$、伴随 $8$。
  \item 见第 \ref{sec:prequant} 讲第 5 题。
  \item $\omega_\lambda(X,Y)=\frac12\langle\lambda,[X,Y]\rangle$ 的核是 $\mathfrak t$（$\lambda$ 正则）；
        $SU(2)$ 时 $\omega_{j}$ 的积分给 $2j$。
  \item 零向量 $\chi=(E\otimes t^{-1})^{k-2j+1}v$：$N_+\chi=0$ 要求指数 $\ge0$，即 $2j\le k$。
  \item 逐句标注：定理（BWB）/引用 Segal/物理断言（量子群）/未给判据（可积性）。
  \item 大 $k$ 极限：$H^0$ 的维数趋于 $\Inv(\otimes_iR_i)$ 的维数，因为约束消失。
\end{enumerate}

\subsection*{第 \ref{sec:wzw} 讲（WZW）}
\begin{enumerate}[leftmargin=*]
  \item $c=\frac{k\dim G}{k+h^\vee}$：$SU(3),k=1$ 给 $c=\frac{8}{4}=2$。\hfill[答 $c=2$]
  \item $k=2$：$h_0=0$、$h_{1/2}=\frac{3/4}{4}=\frac{3}{16}$? 不对，用 $h_j=\frac{j(j+1)}{k+2}$：
        $h_{1/2}=\frac{3/4}{4}=\frac{3}{16}$ 与第 \ref{sec:blocks} 讲的 $3/8$ 差因子 2——
        因为那里 $h$ 用的是 $j(j+1)/(k+2)$ 的另一种归一（$2h$ 约定）。比较前先对齐（附录 \ref{app:conventions}）。
  \item 由 OPE 推 $c$：算 $\langle T(z)T(w)\rangle$ 的 $c/2$ 项，用 Sugawara 的系数。
  \item $n=2$ 二点函数：KZ 给幂律，指数 $2h_j$；$h_j$ 由仿射代数定。
  \item $C_2(j\otimes j)=2C_2(j)+2\Omega$：用 Casimir 的可加性与 $\Omega$ 的本征值。
  \item 对：本文 §3.2 的射影平坦联络在 $g=0$ 的显式形式就是 KZ。
  \item 错：$2j\le k$ 来自仿射代数的可积表示，不是 BWB（BWB 对所有支配整权成立）。
\end{enumerate}

\subsection*{第 \ref{sec:surgery} 讲（手术与 S/T）}
\begin{enumerate}[leftmargin=*]
  \item $k=2$ 的 $3\times3$ $S$（指标 $0,\frac12,1$）：
        \hfill[答]
        \[
        S=\frac{1}{\sqrt2}\begin{pmatrix}
        \frac{1}{\sqrt2} & 1 & \frac{1}{\sqrt2}\\[2pt]
        1 & 0 & -1\\[2pt]
        \frac{1}{\sqrt2} & -1 & \frac{1}{\sqrt2}
        \end{pmatrix}\cdot\frac{1}{\sqrt2}
        \]
        自查：$S^2=\mathbf 1$（对称且实）。
  \item $k=1$：$Z(S^3)=S_{00}=\sqrt{\frac{2}{k+2}}\sin\frac{\pi}{k+2}=\frac{1}{\sqrt3}\sin\frac{\pi}{3}=\frac12$。
        \hfill[答 $1/2$]
  \item 由 (4.40) 展开：$T^m$ 作用在 bra 上给 $q^{-1}$，故相位是 $e^{-2\pi\ii(n+m)c/24}$；
        与原文 (4.41) 的 $e^{2\pi\ii(n-m)c/24}$ 差一个 $e^{-2\pi\ii m c/12}$，属约定差（见附录 \ref{app:errata}）。
  \item (4.47) $\to$ (4.50)：用 $S$ 的对称性与融合矩阵的正交性；(4.48) $\to$ (4.51)：把两边除以 $S_{0i}$。
  \item 秩一性 (4.49)：剪开 $L(R_i;R_j,R_k)$ 后只有一个通道，故矩阵秩一。
  \item 错：(4.28) 是手术律（几何），(4.51) 是融合规则的对角化（代数），两者由 (4.36) 的定义桥接。
  \item 错：标架反常在一般手术下会破坏整体相位；只在“标准标架 + 比值“时自消（p.390 的说明）。
\end{enumerate}

\subsection*{第 \ref{sec:verlinde} 讲（Verlinde）}
\begin{enumerate}[leftmargin=*]
  \item $k=2$：$N_{\frac12\frac12 0}=1$、$N_{\frac12\frac12 1}=1$、$N_{\frac12\frac12\frac12}=0$。
        \hfill[答 $(1,1,0)$]
  \item $k=3$：$N_{112}=1$、$N_{222}=1$、$N_{133}=0$（$j=3/2$ 不可积）。
  \item 取 $i=0$：$S_{0l}=S_{l0}$，$\sum_k S_{0k}N_{kl}{}^m=S_{0l}S_{0}{}^m/S_{00}$ 即 (4.46)。
        不能反推，因为 (4.51) 对一般 $i$ 成立，$i=0$ 只是特例。
  \item 两个方向：$S$ 幺正给逆，$S$ 对称给转置；缺一不可。
  \item (4.52)$\to$(4.54)：把 $v^m$ 的指标降下来（$v_m=S_{0m}v^m$ 的逆），否则求和指标不匹配。
  \item 对“物理推导“成立、“数学证明“不成立：p.391 明说 Moore--Seiberg 已证，他给新证明；
        数学证明是 TUY/Faltings。
  \item 错：十步链只用手术律、(4.36) 定义、秩一性三处输入，共形块动力学不进入。
  \item $k=1$：$j=1/2$ 可积但 $j=1$ 不可积，$N_{\frac12\frac12\frac12}=0$ 说明截断；
        $k=2$：同上 $(1,1,0)$。推理只在“截断后“成立。
\end{enumerate}

\subsection*{第 \ref{sec:duality} 讲（两种对偶）}
\begin{enumerate}[leftmargin=*]
  \item $Z(S^2\times S^1)=1$：空理论配分函数；不能读成边界 CFT 配分函数，因为没有边界自由度。
  \item 有能隙 $\Delta$：$Z$ 会带 $e^{-\Delta}$ 修正，破坏拓扑不变性；故拓扑性排除能隙。
  \item $[G]=L$ 给 $c=3l/2G$ 无量纲；量纲分析只能定 $\alpha$，定不出 $3/2$——后者来自 Brown--Henneaux 的中心扩张计算。
  \item 三条反驳：度规（本文无）、边界条件（本文无渐近 AdS）、字典（无体场-边界算子对应）。
  \item 错在：非紧群、经典层面、边界条件三个关键词各缺一；改写不是等同。
  \item 数值相等不构成等同：$c$ 相同只是必要条件；还要动力学与谱。
  \item $Z(S^3)=S_{00}$ 是 $\cH_{T^2}$ 上的一个数（配分函数值），不是边界 CFT 的配分函数。
  \item 环面：$\dim\cH_{T^2}=k+1$ 有限 vs 二维 CFT 的无穷维；等同只到共形块层。
\end{enumerate}

\subsection*{第 \ref{sec:categorification} 讲（范畴化）}
\begin{enumerate}[leftmargin=*]
  \item 三叶结：$J=-q^{-4}+q^{-3}+q^{-1}$；Khovanov 同调的分次维数表查文献，
        按 (11.1) 取交错和应回到该多项式（注意 $q$ 的幂次归一）。
  \item 双分次更强：两分量链环的 $J$ 相同但 Khovanov 同调不同（如 Hopf 链环的两支）。
  \item $\cH_\Sigma$ 取值 $\mathrm{Vect}$：本文的量子化只给向量空间；链复形需要额外的微分与分次，
        即“高一维“的母理论。
\end{enumerate}
